\chapter{Objetivos}
\label{chap:objetivos}

\drop{U}{na} vez analizado el contexto teórico, el marco normativo y las
herramientas existentes en el capítulo~\ref{chap:antecedentes}, es posible
delimitar con precisión los objetivos que guían el desarrollo de \dataqual.
Este capítulo define, en primer lugar, el objetivo general del \ac{TFG} y, a
continuación, lo descompone en cinco objetivos específicos que abordan las
principales dimensiones funcionales de la plataforma.  Se describen, además,
los objetivos transversales de calidad del software que condicionan las
decisiones arquitectónicas y, por último, se establece el alcance del proyecto
junto con sus limitaciones.


%% =========================================================================
\section{Objetivo general}
\label{sec:objetivo-general}
%% =========================================================================

El objetivo general de este \ac{TFG} consiste en \textbf{diseñar, implementar
y validar una plataforma web interactiva y una \ac{API} \ac{REST}ful que
permitan evaluar de forma automatizada la calidad de los datos empleados en
proyectos de \ac{IA}, \ac{ML} y minería de datos}, proporcionando métricas
cuantitativas, visualizaciones interactivas y mecanismos de diagnóstico que
faciliten la identificación y corrección de problemas de calidad en los
conjuntos de datos.

Este objetivo responde a la necesidad, identificada en la
sección~\ref{sec:motivacion}, de disponer de herramientas que superen las
limitaciones de los enfoques ad~hoc~\textemdash \emph{scripts} puntuales,
consultas \ac{SQL} manuales o inspecciones en cuadernos
Jupyter~\textemdash y que, al mismo tiempo, no impongan la barrera técnica
que presentan las bibliotecas basadas exclusivamente en
código~(sección~\ref{sec:estado-arte}).  \dataqual aspira a
ofrecer una solución \emph{full-stack}, alineada con los estándares
\ac{ISO}/IEC~25012~\cite{iso25012} e \ac{ISO}/IEC~5259~\cite{iso5259}, que
democratice la evaluación de calidad de datos haciéndola accesible tanto a
perfiles técnicos como a usuarios de negocio.


%% =========================================================================
\section{Objetivos específicos}
\label{sec:objetivos-especificos}
%% =========================================================================

El objetivo general se descompone en los siguientes cinco objetivos específicos
(\textbf{OE1}--\textbf{OE5}), cada uno de los cuales aborda una dimensión
funcional diferenciada de la plataforma.


%% -------------------------------------------------------------------------
\subsection{OE1.\ Módulo de ingesta y almacenamiento}
\label{sec:oe1}
%% -------------------------------------------------------------------------

\textbf{Enunciado.}  Desarrollar un módulo de ingesta y almacenamiento que
permita la carga, validación de formato y persistencia segura de conjuntos de
datos.

Este objetivo abarca los siguientes requisitos funcionales:

\begin{itemize}
  \item \textbf{Soporte multiformato.}  La plataforma debe aceptar conjuntos
    de datos en los formatos más utilizados en proyectos de \ac{IA}:
    \ac{CSV}, \ac{XLSX}, \ac{JSON} y Parquet.

  \item \textbf{Parseo robusto.}  Para archivos \ac{CSV}, el módulo debe
    implementar una estrategia de \emph{fallback} que combine múltiples
    codificaciones (UTF-8, UTF-8-BOM, CP1252, Latin-1), separadores (coma,
    punto y coma, tabulador, \emph{pipe}) y motores de lectura, garantizando
    la ingesta correcta del mayor número posible de archivos reales.

  \item \textbf{Almacenamiento seguro.}  Los archivos originales deben
    almacenarse en un sistema de almacenamiento de objetos compatible con S3
    (MinIO), mientras que los metadatos~\textemdash esquema, estadísticas
    descriptivas, conteos~\textemdash se persisten en PostgreSQL.

  \item \textbf{Extracción automática de esquema.}  El sistema debe inferir
    automáticamente los nombres de columna, los tipos de dato y las
    estadísticas descriptivas básicas (media, mediana, desviación estándar,
    mínimo, máximo) de cada conjunto de datos cargado.

  \item \textbf{Versionado de conjuntos de datos.}  El módulo debe soportar
    la carga de nuevas versiones de un mismo conjunto de datos, manteniendo
    la trazabilidad entre versiones mediante una relación padre-hijo.
\end{itemize}


%% -------------------------------------------------------------------------
\subsection{OE2.\ Motor de evaluación con métricas parametrizables}
\label{sec:oe2}
%% -------------------------------------------------------------------------

\textbf{Enunciado.}  Implementar un motor de evaluación que calcule métricas
de calidad estándar sobre los conjuntos de datos cargados, permitiendo su
parametrización según la tipología del proyecto.

Las métricas implementadas deben cubrir las principales dimensiones de calidad
definidas en \ac{ISO}/IEC~25012~\cite{iso25012} y
\ac{ISO}/IEC~25024~\cite{iso25024}, complementadas con métricas específicas
para el contexto de \ac{IA} y \ac{ML}.  En concreto, el motor debe incluir
las siguientes siete métricas:

\begin{enumerate}
  \item \textbf{Completitud} (\texttt{completeness}): ratio de valores no
    nulos por columna, alineada con la característica de completitud de
    \ac{ISO}/IEC~25012.

  \item \textbf{Exactitud sintáctica} (\texttt{syntactic\_accuracy}):
    conformidad de los valores con patrones de formato definidos mediante
    expresiones regulares, correspondiente a la dimensión de exactitud.

  \item \textbf{Consistencia lógica} (\texttt{logical\_consistency}):
    verificación de reglas condicionales entre columnas (reglas
    \texttt{IF\,--\,THEN}), alineada con la dimensión de consistencia.

  \item \textbf{Actualidad} (\texttt{currentness}): antigüedad del dato más
    reciente respecto a un instante de referencia, correspondiente a la
    dimensión de actualidad (\emph{currentness}).

  \item \textbf{Unicidad} (\texttt{uniqueness}): detección de registros
    duplicados, relacionada con la credibilidad de los datos.

  \item \textbf{Valores atípicos} (\texttt{outliers}): detección de
    \emph{outliers} mediante métodos estadísticos (\ac{IQR} y Z-score),
    métrica complementaria no contemplada directamente en
    \ac{ISO}/IEC~25012 pero esencial en el contexto de \ac{ML}.

  \item \textbf{Equilibrio de clases} (\texttt{class\_balance}): evaluación
    del balance de la variable objetivo mediante la entropía de Shannon,
    métrica específica para proyectos de \ac{ML} que permite detectar
    desequilibrios que sesgan los modelos~\cite{ng2021datacentric}.
\end{enumerate}

El motor debe adoptar una arquitectura de \emph{plugins}, de modo que cada
métrica se implemente como una clase independiente que herede de una clase
base abstracta y se registre dinámicamente en un registro centralizado.
Esta arquitectura garantiza la extensibilidad del sistema, permitiendo
incorporar nuevas métricas sin modificar el núcleo del motor.

Adicionalmente, cada métrica debe ser parametrizable: umbrales de aceptación,
columnas objetivo, métodos de detección y pesos relativos (en una escala de
0,0 a 1,0) que determinen su contribución a la puntuación global de calidad.
El sistema debe soportar también plantillas de métricas que permitan
reutilizar configuraciones entre proyectos.


%% -------------------------------------------------------------------------
\subsection{OE3.\ \acs{API} \acs{REST}ful documentada}
\label{sec:oe3}
%% -------------------------------------------------------------------------

\textbf{Enunciado.}  Diseñar e implementar una \ac{API} \ac{REST}ful que
exponga \emph{endpoints} para la carga de datos, configuración de
evaluaciones, ejecución de análisis y consulta de resultados, permitiendo la
integración con flujos de trabajo externos.

La \ac{API} debe organizarse en módulos funcionales cohesivos, siguiendo el
patrón \emph{Blueprint} de Flask:

\begin{definitionlist}
  \item[\texttt{auth}] Registro de usuarios, inicio de sesión, refresco de
    \emph{tokens} \ac{JWT} y cierre de sesión.

  \item[\texttt{projects}] Operaciones \ac{CRUD} sobre proyectos, incluyendo
    la asignación de conjuntos de datos y la configuración de métricas.

  \item[\texttt{datasets}] Carga de conjuntos de datos, consulta de
    metadatos, gestión de versiones y descarga de archivos.

  \item[\texttt{metrics}] Definición y configuración de métricas de calidad,
    gestión de plantillas y consulta del catálogo de métricas disponibles.

  \item[\texttt{evaluations}] Lanzamiento de evaluaciones, consulta de
    progreso, recuperación de resultados e \emph{issues} detectados.

  \item[\texttt{admin}] Utilidades de administración del sistema, incluyendo
    un \emph{endpoint} de comprobación de salud (\emph{health check}).
\end{definitionlist}

Todos los \emph{endpoints} protegidos deben requerir autenticación mediante
\ac{JWT}.  El formato de respuesta debe ser uniforme y predecible, siguiendo
la estructura \texttt{\{success, data, message\}}, con soporte de paginación
en los listados.


%% -------------------------------------------------------------------------
\subsection{OE4.\ Aplicación web con \emph{dashboards} interactivos}
\label{sec:oe4}
%% -------------------------------------------------------------------------

\textbf{Enunciado.}  Desarrollar una aplicación web para la gestión de
proyectos y conjuntos de datos, que permita configurar evaluaciones y umbrales
de calidad personalizados, consultar el historial de evaluaciones realizadas
y visualizar las métricas obtenidas mediante cuadros de mando interactivos.

La interfaz de usuario debe proporcionar las siguientes capacidades:

\begin{itemize}
  \item \textbf{Cuadro de mando principal} con resumen de proyectos,
    conjuntos de datos y distribución de \emph{Quality Gates}.

  \item \textbf{Página de detalle de proyecto} con pestañas diferenciadas
    para conjuntos de datos, configuración de métricas, historial de
    evaluaciones y estado del \emph{Quality Gate}.

  \item \textbf{Visualización de resultados de evaluación} con indicadores de
    puntuación de calidad (\emph{gauge}), pestañas por métrica y listado de
    \emph{issues} detectados.

  \item \textbf{Exploración interactiva de datos} (\ac{EDA}): histogramas,
    diagramas de caja, gráficos de barras, matrices de correlación (Pearson
    y Spearman) y detección visual de valores atípicos con factor \ac{IQR}
    configurable.

  \item \textbf{Gráficos de tendencia} que muestren la evolución de la
    calidad a lo largo de las sucesivas evaluaciones.

  \item \textbf{Indicadores visuales de \emph{Quality Gate}} con códigos de
    color: verde (\texttt{PASSED}), amarillo (\texttt{WARNING}) y rojo
    (\texttt{FAILED}).

  \item \textbf{Diseño adaptable} (\emph{responsive}) con soporte para
    dispositivos de escritorio, tableta y teléfono móvil.
\end{itemize}


%% -------------------------------------------------------------------------
\subsection{OE5.\ Mecanismos de identificación de problemas}
\label{sec:oe5}
%% -------------------------------------------------------------------------

\textbf{Enunciado.}  Desarrollar mecanismos de identificación y diagnóstico
de problemas de calidad que permitan localizar registros defectuosos en los
conjuntos de datos, clasificar las incidencias por severidad y rastrear su
evolución entre evaluaciones sucesivas.

Este objetivo se materializa en tres componentes principales:

\begin{enumerate}
  \item \textbf{Sistema de \emph{issues}.}  Cada problema de calidad
    detectado debe registrarse como un \emph{issue} con cuatro niveles de
    severidad (\emph{critical}, \emph{high}, \emph{medium}, \emph{low}),
    calculada dinámicamente en función de la distancia al umbral configurado.
    Cada \emph{issue} debe incluir muestras de los datos afectados (hasta
    cinco filas de ejemplo), con enmascaramiento automático de columnas
    marcadas como \ac{PII}.

  \item \textbf{\emph{Fingerprinting} determinista.}  Cada \emph{issue} debe
    recibir un identificador único generado mediante \ac{SHA}-256 a partir de
    sus atributos característicos (métrica, columna, tipo de problema),
    permitiendo su identificación unívoca entre ejecuciones independientes.
    Este mecanismo, inspirado en el sistema de rastreo de \emph{issues} de
    SonarQube~\cite{sonarqube}, permite clasificar automáticamente cada
    incidencia como \emph{nueva}, \emph{recurrente} o \emph{corregida}
    respecto a la evaluación anterior.

  \item \textbf{\emph{Quality Gates}.}  Umbrales mínimos configurables que
    determinan si un conjunto de datos es apto para su uso, con tres estados
    posibles: \texttt{PASSED} (todos los umbrales superados),
    \texttt{WARNING} (margen estrecho respecto a algún umbral) y
    \texttt{FAILED} (al menos un umbral incumplido).  Las \emph{Quality
    Gates} proporcionan una señal clara y automatizable que puede integrarse
    en \emph{pipelines} de \ac{MLOps} o de integración continua para
    bloquear el uso de conjuntos de datos de calidad insuficiente.
\end{enumerate}


%% =========================================================================
\section{Objetivos transversales}
\label{sec:objetivos-transversales}
%% =========================================================================

Además de los objetivos funcionales descritos, el desarrollo de \dataqual
persigue un conjunto de objetivos transversales de calidad del software que
condicionan las decisiones arquitectónicas y tecnológicas del proyecto:

\begin{definitionlist}
  \item[Seguridad]  Autenticación basada en \ac{JWT} con \emph{tokens} de
    acceso y refresco, \emph{hashing} de contraseñas mediante \ac{PBKDF2},
    limitación de tasa de peticiones (\emph{rate limiting}), validación
    exhaustiva de las entradas del usuario, protección contra inyección de
    código en las reglas de consistencia lógica y enmascaramiento automático
    de columnas marcadas como \ac{PII}.

  \item[Escalabilidad]  Procesamiento asíncrono de evaluaciones mediante
    Celery y Redis como \emph{broker} de mensajes, lo que permite escalar
    horizontalmente el número de \emph{workers} en función de la carga de
    trabajo, así como separación de responsabilidades entre servicios.

  \item[Mantenibilidad]  Arquitectura modular basada en \emph{blueprints},
    servicios y modelos, con un sistema de métricas extensible por
    \emph{plugins} que permite incorporar nuevas métricas con un esfuerzo
    mínimo.  Migraciones de base de datos versionadas y separación estricta
    entre \emph{frontend} y \emph{backend}.

  \item[Usabilidad]  Interfaz de usuario basada en Material Design con
    formularios dotados de validación en tiempo real, notificaciones
    \emph{toast}, migas de pan (\emph{breadcrumbs}) de navegación y estados
    de carga y error claramente diferenciados.

  \item[Portabilidad]  Despliegue del \emph{stack} completo mediante Docker
    Compose, con configuración por variables de entorno y sin dependencias de
    servicios en la nube propietarios.  Los ocho servicios que componen la
    plataforma~\textemdash \emph{frontend}, \emph{backend}, PostgreSQL,
    MinIO, Redis, Celery \emph{worker}, Celery Beat y
    Flower~\textemdash se orquestan de forma conjunta, garantizando la
    reproducibilidad del entorno y la soberanía sobre los datos.
\end{definitionlist}


%% =========================================================================
\section{Alcance y limitaciones}
\label{sec:alcance}
%% =========================================================================

Para delimitar con claridad el ámbito del presente \ac{TFG}, resulta necesario
establecer tanto lo que forma parte del alcance del proyecto como aquello que
queda excluido.

\subsection{Dentro del alcance}

\begin{itemize}
  \item Diseño e implementación de la plataforma \dataqual como aplicación
    web \emph{full-stack} con \emph{frontend} en Next.js y \emph{backend} en
    Flask.

  \item Implementación de las siete métricas de calidad descritas en el
    objetivo~OE2 (sección~\ref{sec:oe2}), con parametrización completa y
    sistema de pesos.

  \item Desarrollo de la \ac{API} \ac{REST}ful con los seis módulos
    definidos en el objetivo~OE3 (sección~\ref{sec:oe3}).

  \item Despliegue contenerizado del \emph{stack} completo mediante Docker
    Compose.

  \item Validación de la plataforma con conjuntos de datos reales,
    verificando el correcto funcionamiento de las métricas, las
    \emph{Quality Gates} y el sistema de \emph{fingerprinting}.

  \item Documentación técnica: especificación de la \ac{API} y catálogo de
    métricas (anexos del presente documento).
\end{itemize}

\subsection{Fuera del alcance}

\begin{itemize}
  \item \textbf{Procesamiento de \emph{big data}.}  La plataforma está
    diseñada para conjuntos de datos de tamaño moderado (hasta 100\,MB por
    archivo).  La evaluación de conjuntos de datos masivos que requieran
    procesamiento distribuido (Apache Spark, Dask) queda fuera del alcance.

  \item \textbf{Evaluación en tiempo real} (\emph{streaming}).  \dataqual
    opera sobre conjuntos de datos estáticos cargados por el usuario.  La
    evaluación continua de flujos de datos en tiempo real no forma parte del
    alcance actual.

  \item \textbf{Corrección automática de datos.}  La plataforma identifica y
    diagnostica problemas de calidad, pero no implementa mecanismos de
    corrección automática de los datos (limpieza, imputación, deduplicación).

  \item \textbf{Integración nativa con plataformas en la nube.}  Aunque la
    arquitectura basada en MinIO es compatible con S3, no se desarrollan
    conectores específicos para servicios como AWS, Google Cloud o Azure.

  \item \textbf{Internacionalización.}  La interfaz de usuario se desarrolla
    únicamente en español, sin soporte multiidioma.
\end{itemize}

\subsection{Restricciones tecnológicas}

El desarrollo del proyecto está sujeto a las siguientes restricciones
tecnológicas, derivadas tanto de requisitos académicos como de decisiones
arquitectónicas:

\begin{itemize}
  \item El \emph{backend} se implementa en Python con el \emph{framework}
    Flask, utilizando PostgreSQL como sistema de gestión de bases de datos
    relacional.

  \item El \emph{frontend} se desarrolla con React y Next.js, utilizando
    TypeScript como lenguaje principal.

  \item El procesamiento asíncrono se delega en Celery con Redis como
    \emph{broker} de mensajes.

  \item Todo el sistema debe poder desplegarse mediante Docker Compose en
    una única máquina, sin requerir infraestructura distribuida.

  \item El proyecto debe alinearse con los estándares
    \ac{ISO}/IEC~25012~\cite{iso25012},
    \ac{ISO}/IEC~25024~\cite{iso25024} e
    \ac{ISO}/IEC~5259~\cite{iso5259} en lo relativo a las dimensiones
    y métricas de calidad de datos implementadas.
\end{itemize}


% Local Variables:
%  coding: utf-8
%  mode: latex
%  mode: flyspell
%  ispell-local-dictionary: "castellano8"
% End:
